\documentclass[conference]{IEEEtran}

\usepackage[utf8]{inputenc}
\usepackage[T1]{fontenc}
\usepackage{cite}
\usepackage{amsmath}
\usepackage{graphicx}
\usepackage{booktabs}
\usepackage{url}
\usepackage[hidelinks]{hyperref}

\begin{document}

\title{2-Way Set-Associative Cache with LRU Replacement:\\Phase~II Comparative Evaluation}

\author{
\IEEEauthorblockN{Krishna Pramod Dubale, Apoorva Choudhary, Prachee Dhar, Kushal Sarkar, Vani Rudra}
\IEEEauthorblockA{Modern Computer Architecture (MCA) --- Project~7\\
Rishihood University}
}

\maketitle

\begin{abstract}
We implemented a two-way set-associative cache with one-bit
least-recently-used (LRU) replacement and compared it with a direct-mapped
cache. Both Verilog designs store 128 bytes in 16-byte blocks and support
aligned 32-bit reads. We replayed five traces, totaling 89 requests per
design, and ran six simulation suites with 20,193 passing checks. The
two-way cache hit on 14 of 16 requests in an alternating two-block trace,
while the direct-mapped cache missed on all 16. On a three-block cycle,
direct mapping hit on 9 of 30 requests and the two-way cache missed on all
30. The sequential and spatial-locality traces tied. Mean response
latencies over the five traces were 2.506 and 2.584 simulation cycles,
respectively. Routed ECP5 estimates were 102.64 and 137.55~MHz.
\end{abstract}

\begin{IEEEkeywords}
cache memory, set associativity, LRU replacement, direct-mapped cache,
conflict miss, Verilog, hardware verification
\end{IEEEkeywords}

\section{Introduction}

Two frequently accessed blocks can evict each other if a direct-mapped
cache assigns them the same slot. A second way can keep both blocks in the
cache. It also changes how many sets fit in the same capacity, so a
different address sequence can favor direct mapping. Phase~I demonstrated
the first case. Phase~II adds longer traces, a case that favors the
direct-mapped design, and tests that check replacement and request timing
independently.

\section{Literature Review}

Smith~\cite{smith1982} surveys the effects of block size and placement on
cache behavior. Hill~\cite{hill1988} examines the trade-off between miss
rate and access time for direct-mapped caches. We keep data capacity and
block size equal to compare placement on the same addresses.
Jouppi~\cite{jouppi1990} addresses conflict misses with a separate victim cache.
Our design uses two ways in each set and tracks recency with one bit per
set.

Average memory access time (AMAT) relates hit time and miss
penalty~\cite{patterson2017}. We use that model to show the effect of an
assumed refill penalty. The simulation also measures response cycles
directly with the project's combinational memory model.

\section{Proposed Solution}

\subsection{Design Overview}

Each cache has eight lines. The two-way design groups them into four sets
of two ways, and the direct-mapped design uses eight one-way sets.
Table~\ref{tab:address} shows the address fields. Bits [1:0] are zero for
aligned reads.

\begin{table}[ht]
\centering
\caption{Fields of the 32-bit Byte Address}
\label{tab:address}
\begin{tabular}{@{}lcccc@{}}
\toprule
Design & Tag & Set & Word & Byte \\
\midrule
Two-way       & [31:6] & [5:4] & [3:2] & [1:0] \\
Direct-mapped & [31:7] & [6:4] & [3:2] & [1:0] \\
\bottomrule
\end{tabular}
\end{table}

\subsection{Architecture and Controller}

The top-level controller accepts one request at a time on a rising edge
with both \texttt{req\_valid} and \texttt{req\_ready} high. It saves the
address in IDLE and checks the tag in LOOKUP. The two-way storage module
compares both valid tags in parallel and selects a word from the matching
way. A hit goes straight to RESP. A miss goes through REFILL, which writes
the 128-bit block and takes the requested word from the incoming block.
Figure~\ref{fig:datapath} shows the controller and datapath.
Both paths give a one-cycle \texttt{resp\_valid} pulse in RESP before
returning to IDLE. Synchronous reset invalidates all lines.

\begin{figure}[ht]
\centering
\footnotesize
\begin{tabular}{c}
\fbox{CPU read request}\\
$\downarrow$\\
\fbox{Controller: IDLE, LOOKUP, REFILL, RESP}\\
$\updownarrow$ \hspace{3.2em} $\updownarrow$\\[-1pt]
\fbox{Way 0 / Way 1, tag checks, LRU}\quad\fbox{Refill port}\\
$\downarrow$\\
\fbox{Word selection and response}
\end{tabular}
\caption{Two-way cache datapath. A miss refills the selected way and
returns one word from the incoming block.}
\label{fig:datapath}
\end{figure}

\subsection{LRU Replacement}

A refill takes the first invalid way, or the way named by the LRU bit
when both are valid. A hit or refill points the bit to the other way.
In the demo, \texttt{0x00} and \texttt{0x80} fill ways~0 and~1. A hit
on \texttt{0x00} makes way~1 the next victim, so \texttt{0x100} replaces
\texttt{0x80} and the final read of \texttt{0x00} hits.

\section{Results and Discussion}

\subsection{Verification}

The simulation uses Icarus Verilog 13.0~\cite{icarus-verilog}. The memory
model returns four words per block, each equal to its own byte address.
The tests check returned data as well as hit and miss outcomes. The
reference model maintains its own recency order for each set. The
controller tests add reset during an in-flight request, a request held
while the cache is busy, and checks that each completed transaction has
one response.

\begin{table}[ht]
\centering
\caption{Simulation Checks by Test Suite}
\label{tab:verification}
\begin{tabular}{@{}lr@{}}
\toprule
Suite & Passed \\
\midrule
Storage unit & 51 \\
LRU reference-model stress & 19,116 \\
Two-way controller integration & 276 \\
Direct-mapped baseline & 30 \\
Cache comparison & 376 \\
System regression & 344 \\
\midrule
Total & 20,193 \\
\bottomrule
\end{tabular}
\end{table}

\subsection{Trace Comparison}

Both caches start empty for each trace and receive identical addresses
in the same order. Table~\ref{tab:traces} gives the hit counts and rates. In the
extended conflict trace, two blocks repeatedly map to the same set. The
two-way cache holds both after the first two misses. The direct-mapped
cache replaces its only resident line on every access. Both caches tie
on the two traces that reuse words from the same 16-byte blocks.

\begin{table}[ht]
\centering
\caption{Cache Hits per Read and Hit Rate by Trace}
\label{tab:traces}
\small
\begin{tabular}{@{}lrr@{}}
\toprule
Trace & Two-way & Direct \\
\midrule
Short conflict & 2/5 (40.0\%) & 0/5 (0.0\%) \\
Spatial locality & 4/6 (66.7\%) & 4/6 (66.7\%) \\
Extended conflict & 14/16 (87.5\%) & 0/16 (0.0\%) \\
Sequential & 24/32 (75.0\%) & 24/32 (75.0\%) \\
Three-block cycle & 0/30 (0.0\%) & 9/30 (30.0\%) \\
\midrule
Total & 44/89 (49.44\%) & 37/89 (41.57\%) \\
\bottomrule
\end{tabular}
\end{table}

On the three-block cycle \texttt{0x00}, \texttt{0x40}, \texttt{0x80},
all three blocks map to set~0 in the two-way cache, so each access
evicts the next block needed. In the direct-mapped cache, address bit~6
is part of the set index. It places \texttt{0x40} in set~4, where nine
of its ten accesses hit. The other two addresses keep replacing one
another in set~0.

\subsection{Latency and FPGA Evaluation}

The receiver samples a hit two rising edges after request acceptance and
a miss three edges after acceptance. The internal response registers
are updated one edge earlier. Over the 89 requests the two-way cache
took 223 cycles in total, or 2.506 cycles per request. The direct-mapped
cache took 230 cycles, or 2.584 per request.

For the extended conflict trace, $\mathrm{AMAT}=t_{\mathrm{hit}}+(1-h)M$
gives 4.5 and 22.0 cycles for two-way and direct mapping. We use the
trace hit rate $h$, two-cycle hits and a 20-cycle miss penalty $M$.

Yosys and nextpnr placed and routed both cores out of context on an
ECP5 LFE5U-25F, CABGA381, speed grade~6. Both met the 100~MHz target.
The two-way core used more logic and had the lower estimated maximum
clock frequency. These estimates cover the internal core paths.

\begin{table}[ht]
\centering
\caption{Post-Route FPGA Utilization and Timing}
\label{tab:fpga}
\small
\begin{tabular}{@{}lrr@{}}
\toprule
Metric & Two-way & Direct \\
\midrule
Logic cells & 1,817 & 548 \\
Flip-flops & 82 & 79 \\
Distributed RAM write cells & 78 & 39 \\
Block RAM & 0 & 0 \\
Estimated max. clock (MHz) & 102.64 & 137.55 \\
\bottomrule
\end{tabular}
\end{table}

\section{Individual Contributions}

\begin{enumerate}
\item \textbf{Krishna Pramod Dubale} implemented the two-way storage,
tag comparison, and LRU logic, then added the Phase~II reference-model
stress tests.
\item \textbf{Apoorva Choudhary} implemented the cache controller and
memory model, then extended the integration tests for reset, request
handling, and response timing.
\item \textbf{Prachee Dhar} implemented the direct-mapped baseline,
added file-driven benchmarks and recorded hit counts and latency for
both designs.
\item \textbf{Kushal Sarkar} wrote the interface specification, system
regressions, and waveform demo, integrated the six-suite test runner,
fixed regression exit handling, and ran the FPGA comparison.
\item \textbf{Vani Rudra} prepared the IEEE reports, architecture and
controller figures, and the interview handout.
\end{enumerate}

\section{Conclusion}

The two-way cache reduced misses when two active blocks shared a set.
The three-block cycle showed how the direct-mapped cache's extra index
bit can keep one of those blocks resident. Both designs store the same
amount of data, so the difference in hit count comes from how addresses
are placed and replaced. Both designs met 100~MHz on the ECP5 target. The two-way core used more
logic and had a lower estimated maximum clock rate.

\bibliographystyle{IEEEtran}
\bibliography{references}

\end{document}
